% GENERATED FILE. Edit canonical lessons, then run npm run generate:books.
% canonical-source-hash: fnv1a64:7a424b376fb277b2
% canonical-lessons: FR-C263-aller-retour, FR-C263-rate-le-train, FR-C263-revision-voyage

\chapter{At the Station — un aller-retour, j'ai raté le train}
\label{ch:fr-at-the-station}
\hlchaptermodality{263}

\begin{chapteropening}
\textbf{By the end of this chapter:} \emph{I can buy a return or a one-way ticket, ask when and where the train leaves, and report a problem.}

Complete FR-C263-revision-voyage: buy a ticket, then report a missed train and ask for the next one.
\end{chapteropening}

\section[un aller-retour, un aller simple]{un aller-retour — there and back}

\label{lesson:FR-C263-aller-retour}



\begin{quote}
Say \textquotedblleft{}a ticket\textquotedblright{}. (\emph{Un billet.}) Say \textquotedblleft{}the return\textquotedblright{}. (\emph{Le retour.}) Then ask a stranger for a coffee. (\emph{Je voudrais un café.})

You are at the ticket window, and you know the polite opener already.
\end{quote}

\subsection*{You'll want to know: un aller-retour, un aller simple}
\textbf{un aller-retour} — \textquotedblleft{}a return ticket\textquotedblright{}. \textbf{un aller simple} — \textquotedblleft{}a one-way ticket\textquotedblright{}.

Here \emph{aller} is a noun: \textbf{l'aller} is the outward journey, and \textbf{le retour} is the way back. So an \emph{aller-retour} is \textquotedblleft{}a going and a coming back\textquotedblright{}, and an \emph{aller simple} is the going on its own. Both are masculine, \emph{un}, and the word \emph{billet} is usually left out.

\begin{grammarlens}[title={Grammar Lens: at the window}]
The place you are going to takes \textbf{pour}, \textquotedblleft{}for\textquotedblright{}:

\begin{quote}
\emph{\textbf{Je voudrais un aller-retour pour Paris}, s'il vous plaît.} \emph{— Un aller simple pour Lyon, s'il vous plaît.}
\end{quote}

Then the three questions that get you onto the right train, said in one breath: \emph{À quelle heure part le train ? De quel quai ? C'est combien ?} — \textquotedblleft{}What time does the train leave? From which platform? How much is it?\textquotedblright{}

And the answer you will hear back puts the hour and the platform side by side: \emph{Le train part à dix heures, quai trois.}
\end{grammarlens}

\begin{cousinweb}
\textbf{aller} has three Latin parents, and its infinitive is the one usually traced to Latin \emph{ambulāre}, \textquotedblleft{}to walk\textquotedblright{}. \textbf{retour} is formed from the verb for \textquotedblleft{}to return\textquotedblright{}, \emph{re-} plus \textbf{tourner}, \textquotedblleft{}to turn\textquotedblright{}, which is inherited from Latin \emph{tornāre} (EtymDB 2.1). \textbf{simple} is inherited from Latin \emph{simplex}, \textquotedblleft{}single, plain\textquotedblright{} (EtymDB 2.1). So \emph{un aller simple} is, at root, \textquotedblleft{}a single walk out\textquotedblright{}, and the ticket home is a turning back.
\end{cousinweb}

\subsection*{Your turn}
Say these aloud:

\begin{itemize}
\raggedright
  \item \emph{Je voudrais un aller-retour pour Paris, s'il vous plaît.}
  \item \emph{Un aller simple pour Lyon.}
  \item the three questions — \emph{À quelle heure part le train ? De quel quai ? C'est combien ?}
  \item a question for the timetable — \emph{Excusez-moi, où est l'horaire des trains ?}
\end{itemize}

\begin{tcolorbox}[breakable,colback=teal!4,colframe=teal!35!black,title={Before you move on}]
Which ticket brings you home again? (\emph{\textbf{Un aller-retour.}}) Which word comes before the city? (\emph{\textbf{Pour.}}) What is \emph{l'aller}, as a noun? (\textbf{The outward journey}, the going.)
\end{tcolorbox}
\section[j'ai raté le train]{j'ai raté le train — when it goes wrong}

\label{lesson:FR-C263-rate-le-train}



\begin{quote}
Ask for a one-way ticket to Paris. (\emph{Un aller simple pour Paris, s'il vous plaît.})

Now the train has gone without you.
\end{quote}

\subsection*{You'll want to know: j'ai raté le train}
\textbf{j'ai raté le train} — \textquotedblleft{}I've missed the train\textquotedblright{}.

\textbf{rater} is a regular \emph{-er} verb, the everyday one for anything that leaves without you:

\begin{quote}
\textbf{J'ai raté le train.} — I've missed the train. \textbf{J'ai raté le métro.} — I've missed the metro.
\end{quote}

You also have \textbf{manquer}, which says the same thing a little more carefully: \emph{j'ai manqué le train}. In Québec, \emph{manquer} is the more usual verb for it. And for a thing that is gone, not a train that has left, the verb is \emph{perdre}: \emph{j'ai perdu ma valise}, I've lost my suitcase.

English hesitates between \textquotedblleft{}I missed\textquotedblright{} and \textquotedblleft{}I've missed\textquotedblright{}. French uses the passé composé for both: one finished event.

\begin{grammarlens}[title={Grammar Lens: reporting a problem}]
A complaint to a stranger has three parts, and the first is a courtesy:

\begin{quote}
\textbf{Excusez-moi, j'ai un problème.} — a courtesy. \textbf{J'ai raté le train.} — what went wrong. \textbf{Le prochain train part à quelle heure ?} — what you need.
\end{quote}

Or the fault is not yours: \emph{le train est en retard}, \textquotedblleft{}the train is late\textquotedblright{}. Say what happened, then ask the question that gets you moving again.
\end{grammarlens}

\begin{cousinweb}
\textbf{rater} has no parent in EtymDB 2.1, which lists among its senses a gun that misfires. Where the word itself comes from is not settled by either source, so it is left open here. \textbf{manquer} was borrowed from Italian \emph{mancare}, \textquotedblleft{}to be lacking\textquotedblright{}, which goes back to Latin \emph{mancus}, \textquotedblleft{}maimed, powerless\textquotedblright{} (EtymDB 2.1). \textbf{retard} comes from the verb \emph{retarder}, which goes back to Latin \emph{retardāre}, \textquotedblleft{}to keep back, to delay\textquotedblright{} (Lewis \& Short, under \emph{retardo}).
\end{cousinweb}

\subsection*{Your turn}
Say these aloud:

\begin{itemize}
\raggedright
  \item \emph{J'ai raté le train.} — then \emph{J'ai raté le métro.}
  \item the three parts — \emph{Excusez-moi, j'ai un problème. J'ai raté le train. Le prochain train part à quelle heure ?}
  \item the fault that is not yours — \emph{Le train est en retard.}
  \item what you need, with this ticket — \emph{Je peux prendre le prochain train avec ce billet ?}
\end{itemize}

\begin{tcolorbox}[breakable,colback=teal!4,colframe=teal!35!black,title={Before you move on}]
Say \textquotedblleft{}I've missed the train\textquotedblright{}. (\emph{\textbf{J'ai raté le train.}}) Which verb says it more carefully, and is the more usual one in Québec? (\emph{\textbf{Manquer.}}) Name the three parts of a complaint. (A \textbf{courtesy}, \textbf{what went wrong}, \textbf{what you need}.)
\end{tcolorbox}
\section[(review)]{Review — a journey that goes wrong, and is put right}

\label{lesson:FR-C263-revision-voyage}



\begin{quote}
Say \textquotedblleft{}a return ticket\textquotedblright{}. (\emph{Un aller-retour.}) Then \textquotedblleft{}I've missed the train\textquotedblright{}. (\emph{J'ai raté le train.})
\end{quote}

\begin{grammarlens}[title={What you've built}]
Two chapters, and a whole journey you can handle alone. Find the way: \emph{excusez-moi, pour aller à… ?} Buy the ticket: \emph{un aller-retour pour…}. And when it goes wrong, say so in three parts: a courtesy, what happened, what you need.

Here is the whole of it, told afterwards as a story:

\begin{quote}
\emph{J'avais un aller-retour pour Paris. Quand je suis arrivé sur le quai, le train est parti. Donc j'ai demandé : « Excusez-moi, le prochain train part à quelle heure ? »}
\end{quote}

The scene in the imparfait, the events in the passé composé, and the question quoted whole. \emph{Est parti} is \emph{partir} in the passé composé, with \emph{être}.
\end{grammarlens}

\subsection*{Your turn}
Say these aloud:

\begin{itemize}
\raggedright
  \item at the window — \emph{Je voudrais un aller-retour pour Paris, s'il vous plaît.}
  \item on the platform — \emph{Excusez-moi, j'ai un problème. J'ai raté le train.}
  \item the story, afterwards — \emph{Quand je suis arrivé sur le quai, le train est parti.}
  \item directions to the platform — \emph{Allez tout droit, puis tournez à gauche.}
\end{itemize}

\begin{tcolorbox}[breakable,colback=teal!4,colframe=teal!35!black,title={Before you move on}]
What are the three parts of a complaint? (A \textbf{courtesy}, \textbf{what went wrong}, \textbf{what you need}.) Which ticket is one-way? (\emph{\textbf{Un aller simple.}}) In the story, why is \emph{j'avais} in the imparfait? (The ticket in your pocket is \textbf{scene}.)
\end{tcolorbox}
